\section{Opening quotient cycles symmetrically}\label{sec:cycles}
Make one node for each outgoing unit. Its target pair belongs to another
unit or is rooted, so every nonsink node has one outgoing arc.
A nonself unit has two ports, its representative block and its reflected
block; each is internally connectable. A self unit is internally
connected once its core is joined, so its conceptual ports coincide.
A planned link pair joins corresponding source and target ports,
possibly exchanging them. A quotient path therefore lifts to paths
reaching endpoint ports. A twisted quotient cycle need not lift to two
disjoint cycles, and we do not assume that it does.

\subsection{Selfloops and maximal contacts}\label{sec:contacts}
For a quotient selfloop, source and target have equal period $p$ and
different aligned text, even if the target is the opposite orientation.
Thus $|R|<p$. Omit both planned links, join the unit internally, and
use the saved $2p$ to root $R$ and $\rho(R)$ for $2|R|<2p$.
Both ports are rooted.

For a cycle of length at least two, choose a unit with greatest canonical
forward text, breaking ties by stable IDs. Every forward rotation of
\emph{either} orientation of the target text is at most this chosen
source. Let $q$ be the target period.

\begin{lemma}[Paired short contact]\label{lem:cycle-contact}
The selected link pair can be realized within its reserved budget with
one contact of length at most $q$ and its reflection.
\end{lemma}
\begin{proof}
For an individual plan, equation~\eqref{eq:individual} gives the bound $\le q$.
For a collective plan consider its representative block in the nonself
case, or its odd core in the self case. Write $K=kp$, with record
$R=[a,b)$, distinguished start $t$, and $H=f_H(t)$.
If a constituent already has a vertex of length at most $q$, the
collective construction suffices. Otherwise all constituent vertices
have length greater than $q$.

The maximal-rotation contract and the canonical comparison give
$b-t<q$. Put $x=\min(t,a+K)$. If $x=t$ this already gives $b-x<q$;
otherwise
$b-x=b-a-K<p+q-K\le q$. Thus $[x,b)$ is a short suffix of $R$.
Every constituent first end at $x$ exceeds $b$, since its vertex length
exceeds $q>b-x$, and is at most $H$. With $F=f_H(x)$ we have
$b<F\le H\le b+K$. Follow the ordered list
\begin{equation}\label{eq:cycle-list}
[a,b),\quad[x,b),\quad[x,F),\quad[a+K,b+K).
\end{equation}
Its cost is $K$; the growth at $x$ touches all constituents, and the
short suffix is the desired contact. Its prescribed windows are actual
vertices or substrings of them. This is the geometric proof of
\src{\texttt{lem:cycle-contact}}, now applied to the odd core where
appropriate, not an unqualified use of its original block statement.
The final self-layer meets its core partner. Reflect the constructed
walk to obtain the second contact within the paired budget.
\end{proof}

\begin{theorem}[Paired connection budget]\label{thm:connection}
One can add paired closed walks of total up cost at most $W$ to $G$
so that the resulting graph $D$ is connected to $\eps$. It is balanced,
$\tau$-invariant, and has even outdegree at every fixed vertex.
Its total up cost $U$ is at most $2W$.
\end{theorem}
\begin{proof}
For a cycle of length at least two retain the selected link pair of
Lemma~\ref{lem:cycle-contact}, but omit the \emph{target unit's} outgoing
pair. These are different links. Join that target internally, saving
$2q$, and use the saving to root both short contacts. The quotient
cycle with the omitted edge removed is a path ending at the target.
Every lifted component meets one of the retained contact pieces, which
touch both target ports. Rootedness propagates along the path to every
source port, whether the lift was twisted or untwisted. This also
covers mixed self and nonself cycles.

Different quotient cycles have disjoint accounts. Execute all other
links in increasing distance from rooted sinks or opened cycles.
Each uses its source unit's reserved account. Together with the earlier
rooting and request operations, this roots every layer. The accounts
partition the original layer budgets, whose sum is $W$.

Every addition was a closed walk and its reversed partner, with fresh
paired edge copies. Hence balance, symmetry, and fixed-vertex parity
persist by Section~\ref{sec:base}. All base edges remain, so required vertices
remain used, and total up cost is at most $W+W=2W$.
\end{proof}
